% hw-6-ordering.tex

\newpage
\section{作业 6：序关系}

\begin{problem}[集合划分与精化关系]
  令 $X \neq \emptyset$ 为任意集合。
  令 $Pt(X)$ 表示 $X$ 的所有划分组成的集合。

  \noindent 如下定义 $Pt(X)$ 上的精化 (refinement) 关系 $\le$:
  \[
    \alpha \le \beta \iff
      \forall A \in \alpha\; \exists B \in \beta\; A \subseteq B.
  \]
  在课堂上我们已经证明, ($Pt(X)$, $\le$) 是偏序集。

  \fig{width = 0.50\textwidth}{figs/set-partition-refinement}

  \begin{enumerate}[(1)]
    \item 请证明, 任意 $S \subseteq Pt(X)$ 都有上确界, 并给出上确界的表达式。
    \item 请给出任意 $S \subseteq Pt(X)$ 的下确界的表达式。
  \end{enumerate}
\end{problem}

\begin{remark}
	在精化偏序下, ``更小'' 意味着 ``更细''。因此下确界是公共细化, 上确界是公共粗化。\par
	上下确界的构造分别对应于 "合并（合并为连通分支）" 与 "细化（取交）" 两种直观操作: 上确界通过在必要时合并块来获得最小的公共粗化,
	下确界则通过取各划分块的交得到最大的公共细化。两者都可直接从精化关系的定义出发严格推导。
\end{remark}

\begin{proof}
	\begin{enumerate}[(1)]
		\item[上确界]\
			对每个 $\alpha\in S$ 定义其诱导的等价关系
			\[x\equiv_{\alpha} y \iff \exists A\in\alpha\; (x\in A\land y\in A).
			\]
			令
			\[R\;=\;\bigcup_{\alpha\in S} \equiv_{\alpha},\]
			并令 $\approx$ 为包含 $R$ 的最小等价关系（即 $R$ 的反身-对称-传递闭包）。记
			\[P\;=\;X/\approx.
			\]

			断言 $P=\sup(S)$. 证明如下。

			(i) $P$ 是上界：任取 $\alpha\in S$, 若 $A\in\alpha$ 且 $x,y\in A$，则 $(x,y)\in R$, 因而 $x\approx y$;
			故每个 $A\in\alpha$ 完全包含于某个 $\approx$-类， 即 $\alpha\le P$。

			(ii) $P$ 是最小上界：若 $\beta$ 是任意上界 (即 $\forall\alpha\in S\;\alpha\le\beta$)，
			则对任意 $\alpha\in S$ 有 $\equiv_{\alpha}\subseteq\equiv_{\beta}$，于是
			$R\subseteq\equiv_{\beta}$. 因为 $\equiv_{\beta}$ 是等价关系, 它包含 $R$ 的等价闭包, 即
			$\approx\subseteq\equiv_{\beta}$。因此每个 $\approx$-类都包含于某个 $\beta$-类, 也就是 $P\le\beta$。

			由 (i),(ii) 得 $P$ 为最小上界.

			\emph{直观说明.} 将元素看作点, 若两点在某个划分中同块则连一条边; 上确界的块即该无向图的连通分支——
			把在不同划分中被"间接"连接的块合并起来, 只做必要的合并以保证每个原划分的块被包含于某一合并后的块中。

		\item[下确界]\
			令
			\[\Pi=\Big\{\;\bigcap_{\alpha\in S}A_{\alpha}\;\Big|\;A_{\alpha}\in\alpha\ \forall\alpha\in S,\ \bigcap_{\alpha\in S}A_{\alpha}\neq\emptyset\;\Big\}.
			\]

			断言 $\Pi=\inf(S)$. 证明如下。

			(i) $\Pi$ 为划分：对任意 $x\in X$，对每个 $\alpha\in S$ 取包含 $x$ 的块 $A_{\alpha}\in\alpha$，则
			$x\in\bigcap_{\alpha}A_{\alpha}\in\Pi$，所以 $\Pi$ 覆盖 $X$。若
			$C=\bigcap_{\alpha}A_{\alpha}$ 与 $D=\bigcap_{\alpha}B_{\alpha}$ 都属 $\Pi$ 且 $C\cap D\neq\emptyset$，
			则任选 $x\in C\cap D$ 可得 $A_{\alpha}=B_{\alpha}$ 对所有 $\alpha$, 从而 $C=D$，故 $\Pi$ 的块两两不交。

			(ii) $\Pi$ 为下界：显然对任意 $\alpha\in S$ 与任意 $C=\bigcap_{\gamma}A_{\gamma}\in\Pi$ 有 $C\subseteq A_{\alpha}$,
			因此 $\Pi\le\alpha$ 对所有 $\alpha$ 成立。

			(iii) $\Pi$ 为最大下界：若 $\gamma\in Pt(X)$ 也是下界 (即 $\forall\alpha\in S\;\gamma\le\alpha$)，
			则任取 $G\in\gamma$ 对所有 $\alpha$ 存在 $A_{\alpha}\in\alpha$ 使得 $G\subseteq A_{\alpha}$，故有
			$G\subseteq\bigcap_{\alpha}A_{\alpha}\in\Pi$，因此 $\gamma\le\Pi$。

			由 (i),(ii),(iii) 得 $\Pi$ 为最大下界.

			\emph{直观说明.} 把所有划分叠在一起观察: 下确界的块就是所有划分的交集的非空格子——把每个划分都切得更细,
			只留下在每个划分中都在同一块的点集合。
	\end{enumerate}
\end{proof}
%%%%%%%%%%%%%%%

\begin{problem}[拟序、等价关系、偏序]
  一个自反且传递的二元关系 $R \subseteq X \times X$
  称为 $X$ 上的拟序 (或预序)。
  \footnote{
    The name {\it preorder} is meant to suggest that
    preorders are {\it almost} partial orders,
    but not quite, as they are not necessarily antisymmetric.
  }

  现令 $\preceq\; \subseteq X \times X$ 为拟序。

  \begin{enumerate}[(1)]
    \item 如下定义 $X$ 上的关系 $\sim$:
      \[
        x \sim y \triangleq x \preceq y \land y \preceq x.
      \]
      请证明~\footnote{你在 \textsl{hw4} 中已经做过这个证明了,
      不必重做。可以直接在第二问中使用该结论。}, $\sim$ 是 $X$ 上的等价关系。
      \footnote{Each preorder induces an {\it equivalence relation}
        between elements according to $\sim$,
        which is similar to the {\it antisymmetric} property.}
    \item 如下定义商集 $X/\sim$ 上的关系 $\le$:
      \[
        [x]_{\sim} \le [y]_{\sim} \triangleq x \preceq y.
      \]
      请证明, $\le$ 是偏序关系。
      \footnote{Preorders can be turned into partial orders
        by identifying all elements that are equivalent
        with respect to $\sim$.}
    \item 接下来考虑整数集合 $\Z$ 上的整除关系 $\mid$。
      \begin{itemize}
        \item 请证明, $\mid$ 是一个拟序。
        \item 请给出元素 $x \in \Z$ 的等价类 $[x]_{\sim}$。
        \item 请描述商集 $\Z/\sim$ 上的偏序关系 $\le$。
        \item 请说明, 等价关系 $\sim$ 提供了什么样的``抽象''?
      \end{itemize}
    \item 接下来考察有向图 $G = (V, E)$ 上~\footnote{
        $V$ 是顶点集合, $E \subseteq V \times V$ 是边集合。}
        \mfig{width = 0.90\textwidth}{figs/scc}
      的可达关系 $\leadsto$。
      \footnote{对于 $v, w \in V$, $v \leadsto w$ 表示存在一条从 $v$ 到 $w$ 的路径。}
      \begin{itemize}
        \item 请证明, $\leadsto$ 是一个拟序。
        \item 请给出顶点 $v \in V$ 的等价类 $[v]_{\sim}$。
        \item 请描述商集 $V/\sim$ 上的偏序关系 $\le$。
        \item 请说明, 等价关系 $\sim$ 提供了什么样的``抽象''?
      \end{itemize}
  \end{enumerate}
\end{problem}

\begin{proof}
  \begin{intention}
    先把拟序的 ``双向可比部分'' 提炼成等价关系, 再在商集上压缩掉这些等价元素, 从而得到偏序。后两问分别把这个抽象过程落实到整除关系与可达关系。
  \end{intention}

  \begin{enumerate}[(1)]
    \item 下证 $\sim$ 是 $X$ 上的等价关系。

      \begin{enumerate}[(i)]
        \item \textbf{自反性.} 对任意 $x \in X$, 由于 $\preceq$ 自反,
          \[
            x \preceq x.
          \]
          故
          \[
            x \preceq x \land x \preceq x,
          \]
          即 $x \sim x$。

        \item \textbf{对称性.} 对任意 $x, y \in X$, 若 $x \sim y$, 则
          \[
            x \preceq y \land y \preceq x.
          \]
          交换两个合取项得
          \[
            y \preceq x \land x \preceq y,
          \]
          即 $y \sim x$。

        \item \textbf{传递性.} 对任意 $x, y, z \in X$, 若 $x \sim y$ 且 $y \sim z$, 则
          \[
            x \preceq y, \qquad y \preceq x,
            \qquad y \preceq z, \qquad z \preceq y.
          \]
          由 $\preceq$ 的传递性,
          \[
            x \preceq z,
            \qquad
            z \preceq x.
          \]
          故 $x \sim z$。
      \end{enumerate}

      因而 $\sim$ 是 $X$ 上的等价关系。

    \item 下证 $\le$ 是商集 $X/\sim$ 上的偏序关系。

      先证 \textbf{良定义}。若
      \[
        [x]_{\sim} = [x']_{\sim},
        \qquad
        [y]_{\sim} = [y']_{\sim},
        \qquad
        x \preceq y,
      \]
      则
      \[
        x' \sim x,
        \qquad
        y \sim y'.
      \]
      从而
      \[
        x' \preceq x,
        \qquad
        y \preceq y'.
      \]
      再由传递性,
      \[
        x' \preceq x \preceq y \preceq y',
      \]
      故
      \[
        x' \preceq y',
      \]
      即
      \[
        [x']_{\sim} \le [y']_{\sim}.
      \]
      因而 $\le$ 良定义。

      再证三条偏序公理。

      \begin{enumerate}[(i)]
        \item \textbf{自反性.} 对任意 $[x]_{\sim} \in X/\sim$, 由 $\preceq$ 自反,
          \[
            x \preceq x,
          \]
          故
          \[
            [x]_{\sim} \le [x]_{\sim}.
          \]

        \item \textbf{反对称性.} 若
          \[
            [x]_{\sim} \le [y]_{\sim}
            \qquad \text{且} \qquad
            [y]_{\sim} \le [x]_{\sim},
          \]
          则
          \[
            x \preceq y
            \qquad \text{且} \qquad
            y \preceq x.
          \]
          故 $x \sim y$, 从而
          \[
            [x]_{\sim} = [y]_{\sim}.
          \]

        \item \textbf{传递性.} 若
          \[
            [x]_{\sim} \le [y]_{\sim}
            \qquad \text{且} \qquad
            [y]_{\sim} \le [z]_{\sim},
          \]
          则
          \[
            x \preceq y
            \qquad \text{且} \qquad
            y \preceq z.
          \]
          由 $\preceq$ 的传递性,
          \[
            x \preceq z.
          \]
          故
          \[
            [x]_{\sim} \le [z]_{\sim}.
          \]
      \end{enumerate}

      因而 $\le$ 是偏序关系。

    \item 现考虑 $(\Z, \mid)$。

      \begin{enumerate}[(a)]
        \item \textbf{$\mid$ 是拟序.}

          自反性: 对任意 $x \in \Z$,
          \[
            x = 1 \cdot x,
          \]
          故 $x \mid x$。

          传递性: 若 $x \mid y$ 且 $y \mid z$, 则存在 $m, n \in \Z$ 使得
          \[
            y = mx,
            \qquad
            z = ny.
          \]
          从而
          \[
            z = n(mx) = (nm)x,
          \]
          故 $x \mid z$。

          因而 $\mid$ 是拟序。

        \item \textbf{等价类 $[x]_{\sim}$.}

          由定义,
          \[
            x \sim y
            \iff x \mid y \land y \mid x.
          \]

          若 $x = 0$, 则由 $y \mid 0$ 与 $0 \mid y$ 可得 $y = 0$,
          故
          \[
            [0]_{\sim} = \set{0}.
          \]

          若 $x \neq 0$, 且 $x \sim y$, 则存在 $a, b \in \Z$ 使得
          \[
            y = ax,
            \qquad
            x = by.
          \]
          于是
          \[
            x = b(ax) = (ab)x.
          \]
          由于 $x \neq 0$, 得
          \[
            ab = 1.
          \]
          在整数中这只可能发生于
          \[
            a = b = 1
            \qquad \text{或} \qquad
            a = b = -1.
          \]
          故 $y = x$ 或 $y = -x$。

          反过来, 显然
          \[
            x \mid x,
            \qquad
            x \mid (-x),
            \qquad
            (-x) \mid x.
          \]
          故当 $x \neq 0$ 时,
          \[
            [x]_{\sim} = \set{x, -x}.
          \]

        \item \textbf{商集上的偏序关系.}

          映射
          \[
            \Phi : \Z/\sim \to \set{0} \cup \Z_{>0},
            \qquad
            \Phi([x]_{\sim}) = |x|
          \]
          是双射, 且
          \[
            [x]_{\sim} \le [y]_{\sim}
            \iff x \mid y
            \iff |x| \mid |y|.
          \]
          因而 $\Z/\sim$ 上的偏序, 就是把正负号压缩之后得到的 ``按绝对值整除'' 的偏序。

        \item \textbf{抽象的含义.}

          等价关系 $\sim$ 把只相差一个符号的整数识别为同一个对象:
          \[
            x \sim y \iff |x| = |y|.
          \]
          因而该抽象忘掉了 ``正负'', 只保留 ``大小的整除结构''。
      \end{enumerate}

    \item 现考虑有向图 $G = (V, E)$ 上的可达关系 $\leadsto$。

      \begin{enumerate}[(a)]
        \item \textbf{$\leadsto$ 是拟序.}

          自反性: 对任意 $v \in V$, 存在从 $v$ 到 $v$ 的长度为 $0$ 的路径,
          故
          \[
            v \leadsto v.
          \]

          传递性: 若 $u \leadsto v$ 且 $v \leadsto w$, 则存在从 $u$ 到 $v$ 的路径与从 $v$ 到 $w$ 的路径。将二者首尾连接, 得到从 $u$ 到 $w$ 的路径, 故
          \[
            u \leadsto w.
          \]

          因而 $\leadsto$ 是拟序。

        \item \textbf{等价类 $[v]_{\sim}$.}

          由定义,
          \[
            v \sim w
            \iff v \leadsto w \land w \leadsto v.
          \]
          因而
          \[
            [v]_{\sim} = \set{w \in V \mid v \leadsto w \land w \leadsto v}.
          \]
          这正是顶点 $v$ 所在的强连通分量 (strongly connected component, SCC)。

        \item \textbf{商集上的偏序关系.}

          商集 $V/\sim$ 上的关系由
          \[
            [v]_{\sim} \le [w]_{\sim}
            \iff v \leadsto w
          \]
          给出。由第 (2) 问, 该关系良定义且构成偏序。
          换言之, 这就是把每个 SCC 缩成一个点之后得到的可达偏序;
          其 Hasse 图对应于 SCC condensation graph 的无环骨架。

        \item \textbf{抽象的含义.}

          等价关系 $\sim$ 把 ``彼此可达'' 的顶点识别为同一个对象, 即把每个 SCC 压缩成一个点。该抽象忘掉了 SCC 内部的回路细节, 只保留不同 SCC 之间的可达先后关系。
      \end{enumerate}
  \end{enumerate}
\end{proof}
%%%%%%%%%%%%%%%

\begin{problem}[上确界、下确界]
  令 $X \neq \emptyset$。
  在课堂上, 我们已经证明 $(\pow{X}, \subseteq)$ 是偏序集。

  \noindent 请给出任意 $\emptyset \neq S \subseteq \pow{X}$ 的上确界与下确界的表达式
  (无需证明)。
\end{problem}

\begin{solution}
  对任意 $S \subseteq \pow{X}$,
  \[
    \sup(S) = \bigcup_{A \in S} A,
    \qquad
    \inf(S) = \bigcap_{A \in S} A.
  \]
\end{solution}
%%%%%%%%%%%%%%%

\begin{problem}[偏序]
  令 $X, Y$ 为集合, 定义
  \[
    Fn(X, Y) = \bigcup_{Z \subseteq X} Y^Z.
  \]
  即, $Fn(X, Y)$ 是由 $X$ 的子集到 $Y$ 的所有函数组成的集合。

  \noindent 请证明, $\subseteq$ 是 $Fn(X, Y)$ 上的偏序关系。
\end{problem}

\begin{proof}
  \begin{intention}
    将 $Fn(X, Y)$ 中的函数都看成有序对的集合, 则题目就退化为 ``集合包含关系在一个集合族上构成偏序''。
  \end{intention}

  任取 $f, g, h \in Fn(X, Y)$。由于每个函数本身都是有序对的集合, 故对它们讨论 $\subseteq$ 是有意义的。

  \begin{enumerate}[(1)]
    \item \textbf{自反性.}
      \[
        f \subseteq f.
      \]

    \item \textbf{反对称性.} 若
      \[
        f \subseteq g
        \qquad \text{且} \qquad
        g \subseteq f,
      \]
      则由集合的反对称性,
      \[
        f = g.
      \]

    \item \textbf{传递性.} 若
      \[
        f \subseteq g
        \qquad \text{且} \qquad
        g \subseteq h,
      \]
      则由集合包含关系的传递性,
      \[
        f \subseteq h.
      \]
  \end{enumerate}

  因而 $\subseteq$ 在 $Fn(X, Y)$ 上满足自反性、反对称性与传递性, 故它是 $Fn(X, Y)$ 上的偏序关系。
\end{proof}